\section{Propuesta PSO}
\subsection{Datos del mercado}

Los datos utilizados en este trabajo corresponden a precios
históricos de criptomonedas, los cuales fueron obtenidos mediante la
biblioteca \texttt{yfinance} de Python, que permite acceder a
información financiera desde fuentes públicas.

Se seleccionaron las siguientes criptomonedas: Bitcoin, Ethereum,
Ripple y Litecoin, debido a su relevancia dentro del mercado y
disponibilidad de datos históricos.

El periodo de análisis abarca desde [AÑO] hasta [AÑO], utilizando una
frecuencia [diaria/semanal], lo que permite capturar el
comportamiento del mercado y su volatilidad en el tiempo.

A partir de estos datos, se calcularon los rendimientos logarítmicos
de cada activo, los cuales son utilizados como base para la
construcción y evaluación de los portafolios.
\begin{equation}
  r_t = \ln\left(\frac{P_t}{P_{t-1}}\right)
\end{equation}
Donde \( P_t \) representa el precio del activo en el tiempo \( t \).

El motivo de utilizar rendimientos logarítmicos radica en su
capacidad para manejar mejor las variaciones porcentuales y su
propiedad de aditividad, lo que facilita el análisis de los
portafolios y la comparación entre diferentes activos.
\subsection{Como se maneja la volatilidad}
\subsubsection{Funcion objetivo}
\subsubsection{Modelo CGARCH}
\subsection{Parametros del algoritmo}
\section{Resultados}
\section{Analisis de resultados}
\section{Conclusiones}
